A comparison between two mathematical verifiers can depend on a small number of uncertain answer grades. On a fixed generation pool, each answer contributes its grade times the difference in its probabilities of selection. We use this established linear comparison to ask how to spend a budget when preliminary judgments can be wrong. Given an upper bound on their error count, we derive the smallest attainable worst-case interval width both before provisional grades arrive and after they are available for selective adjudication. We replay 1,390,800 eligible responses from two public archives. On AceMath MATH-500 at eight draws, 437 reliable answer checks establish a 1.21 percentage-point lead, but 13 adverse grade changes can erase it. In the Lightman full-pool comparison, assuming at most five errors in either of two automatic grading rules gives intervals that exclude the archived 5.8-point gap. After calibrating the error bound by probability sampling, all six primary intervals contain zero. At the same adjudication budget, directly sampling signed grading errors gives narrower median intervals in all six settings. The minimax design is useful when an error ceiling can be defended and a guarantee across all compatible grades is required. Estimating a ceiling can be costly when the aim is a statistical comparison.
